% 【本文件写什么】api-v0 契约层，叶子，只写语义不写实现
% - crate 路径：os/components/wateros-task/task-scheduler/scheduler-api/api-v0/src/
% - 列出并说明：pub trait、核心类型、错误枚举、常量
% - 每个公开 API 的前置条件、后置条件、线程/中断上下文限制
% - 与 Linux/用户态约定的对齐点与 deliberate 差异
% - v0 版本当前行为与后续可替换 extension 点
% 【事实来源】
% - 上述 crate 下全部 .rs 源文件
% - docs/exports/public-api/ 中对应符号表，以源码为准
% 【不写什么】
% - impl 内部数据结构、汇编、具体算法步骤

\subsubsection{调度 trait 与触发原因}

\code{ScheduleReason}枚举一次调度决策的触发来源：\code{StartFirst}、\code{Yield}、\code{Tick}、\code{Block(TaskBlockReason)}、\code{Sleep(TaskTick)}、\code{Exit(TaskExitCode)}。实现据此更新当前任务状态并选择下一就绪任务。

\code{SwitchScheduler} trait：\code{prepare\_first\_switch}返回首次切换指针对\code{SwitchPair}；\code{schedule}、\code{schedule\_wait}在普通/等待路径产出\code{Option<SwitchPair>}供 arch \code{\_\_switch}消费。

\code{Scheduler} trait 定义对外最小能力集：\code{init}、\code{spawn\_kernel\_task}、\code{spawn\_user\_task}、等待队列分配/释放、\code{run\_first\_task}、\code{schedule}、\code{block\_current}、\code{wait\_current}及其\code{\_while}/\code{\_timeout}变体、\code{wake\_task}、\code{reap\_exited\_task}等。默认方法实现\code{wait\_current\_while}等条件等待组合。

\subsubsection{就绪队列与注册表抽象}

\code{ReadyQueue}、\code{ReadyTaskSink} trait 抽象就绪队列入队/出队，供 multi-class 与 round-robin 共享骨架。\code{TaskRegistry}封装 TCB 槽位、快照查询、fork/clone/exec 创建路径。\code{WaitQueues}管理显式等待队列、睡眠队列与超时提升。

\code{QueueTarget}、\code{SchedPolicyChangeAction}描述策略变更后是否需立即重调度。\code{SchedulableCheck}由\code{task-api}再导出，供\code{set\_scheduler}校验。

\begin{rustcode}
pub trait SwitchScheduler {
    fn prepare_first_switch(&mut self) -> SwitchPair;
    fn schedule(&mut self, reason: ScheduleReason) -> Option<SwitchPair>;
    fn schedule_wait(
        &mut self, wait_handle: TaskWaitHandle,
        timeout_ticks: Option<TaskTick>,
    ) -> Option<SwitchPair>;
}
\end{rustcode}

\subsubsection{上下文限制}

调度操作须在关中断或调度器自管临界区内调用；\code{wait\_current\_while}要求\code{condition}闭包在持有调度锁时执行，避免丢失唤醒。\code{run\_first\_task}返回\code{!}，永不返回至调用方。
